\selectlanguage{english}
\noindent\textbf{Paths, localization, and the information in an indexing category}

\medskip
\noindent\emph{Commentary and audit produced by GPT-6 Astra (OpenAI, ChatGPT Pro mode).  This credit concerns the present work; it does not certify any subsequent assembly or editorial revision.}

\medskip
These pages occupy the middle of Chapter~VII.  They begin with the end of the combinatorial analysis of path types, establish the ordinary category of fractions by means of path categories, and then examine whether a category containing all shapes of paths has the same homotopy type as a sequentially enlarged model.  Their central concern is a compatibility problem.  Concatenation is completely satisfactory before taking limits: it is functorial, associative and unital.  The sequential models make other calculations accessible, but their transition maps need not respect that composition strictly.  Grothendieck is trying to preserve both advantages in one construction.

The category $\underline{\mathrm{Ty}}$ consists of finite oriented intervals, equipped with their spatial order and their two endpoints.  Its arrows preserve that order, the categorical direction of each traversed edge, and the endpoints.  Concatenation makes it a strict monoidal category.  The opening calculation describes its nerve degree by degree as a free monoid: a composable string of type maps splits uniquely over the edges of its initial interval.  The data attached to one initial edge are successive finite weightings recording how many edges replace it at each stage.  The earlier VII-a commentary explains this calculation through manuscript page~57.  The present audit checks the overlap again and adds the distinction needed for its later use: a face can send an indecomposable generator to a product.  Degreewise freeness is therefore not freeness on a simplicial set.

For an ordinary category $X$, a path of type $\tau$ is a functor $I_\tau\to X$.  Between paths of the same type one keeps natural transformations, and between different types one also keeps pullback along interval maps.  These categories assemble into a category fibred over $\underline{\mathrm{Ty}}$.  Fixing endpoints $a,b$ means taking the strict fibre of the endpoint functor: the components of a natural transformation at both ends must be identities.  This requirement is responsible for several of the correct diagrams in the text.  Dropping it would change the hom-sets obtained after passing to components.

If $\Sigma$ is a wide subcategory of $X$, we allow backward edges only when they lie in $\Sigma$, and write $\Ch_\Sigma(X;a,b)$ for the resulting total fixed-endpoint path category.  Closure of $\Sigma$ under identities and composition is exactly what makes this construction stable under every type map.  A source edge can collapse, or cover a consecutive block of target edges; the two possibilities require precisely these two closure conditions.  No calculus-of-fractions hypothesis is needed at this stage.  The fibres are full subcategories of functor categories, so their internal vertical transformations remain unrestricted.

The first decisive result is
\[
 \operatorname{Hom}_{X[\Sigma^{-1}]}(a,b)
      \cong \pi_0\Ch_\Sigma(X;a,b).
\]
Concatenation induces the composition on the right.  A one-edge path followed by its reversal is connected to a constant path through a ladder with identity endpoint components, and the constant path is connected to the empty path by a type collapse.  Thus every arrow of $\Sigma$ becomes invertible.  For a functor $F:X\to Y$ which inverts $\Sigma$, evaluation of a zigzag in $Y$ is unchanged under either a ladder or a type change.  In a ladder, the equations for its successive squares telescope to
\[
 \eta_n E_F(c)=E_F(c')\eta_0;
\]
the endpoint identities give equality of effects.  The argument uses inverses of the backward horizontal edges but never inverses of the internal vertical components.  It supplies the omitted well-definedness check for the map $\beta$ on manuscript page~68.  The resulting evaluation has the required unique factorization property.  This is the classical localization problem of \citet[Chapter~I, \S1]{GabrielZisman1967}, with the witnesses for its relations retained as categorical data.

Two examples make the theorem tangible.  Inverting one arm of a span creates the fraction across the span and produces a category with a walking-arrow skeleton.  Inverting both arrows of the category $a\rightrightarrows b$ produces a connected groupoid with automorphism group $\mathbf Z$; the nerve of the original category is a circle.  The complete calculations distinguish cancellation of backtracking from the survival of a genuine loop.  They also show why a relative localization need not be a groupoid.

The next task is to establish that concatenation really acts on all morphisms of the total path categories.  A ladder can be glued to a second ladder because their common endpoint components are identities.  A change of type can be glued to a second change of type because the common endpoint has a specified image.  Writing a total morphism as a type map together with a vertical transformation makes the functoriality and strict associativity transparent.  The commentary gives this calculation explicitly.  It also repairs an edition diagram which glues over the two-element set $\{0,1\}$ where the construction requires a single identified endpoint.

Sections~8 and~9 then explain why the particular interval map inducing a pullback ceases to matter after categorical homotopy.  The relevant homotopy relation joins functors by finite zigzags of natural transformations.  Grothendieck proves that every nonempty category of endpoint-preserving interval maps is connected, and asks whether it is contractible.  We give a complete categorical contraction.  After a forced initial collapse when the first signs disagree, the first target run can be filled by the first source edge.  Clamping the remaining vertex images to the end of that run defines an idempotent functor and a natural transformation in the direction determined by the run's sign.  The retract is a smaller interval-mapping category.  Induction contracts it to a point.  This is marked as a development of the present commentary, without a claim of priority.

Precomposition transports these natural transformations to the path categories.  Because the endpoint components remain identities, the homotopy takes place over $X\times X$ and restricts to every fixed-endpoint fibre.  In modern notation, the diagram of fibres becomes constant on parallel arrows after passage to $\Cathot$, the category whose morphisms are categorical-homotopy classes of functors.  This is the bridge to the abstract discussion of a fibration $C\to T$ in \S10.

Here three objects must be distinguished.  There is the total category $C$, the ind-object determined by its fibres in $\Cathot$, and an ordinary sequential colimit $C_\infty$ formed after choosing a chain cofinal in the reachability preorder $T_{\mathrm o}$.  The source assumes that this preorder is cofiltered.  It does not assume that parallel arrows of $T$ can be equalized.  Forgetting that difference would conceal the principal obstruction: for the constant point-valued diagram over a one-object group $BG$, the ind-object and every sequential model are points, while the total category is $BG$.

There is a further failure on manuscript page~103.  The reachability preorder of a slice $T/t_0$ need not be the lower ideal of $t_0$ in $T_{\mathrm o}$.  The category with two parallel arrows already disproves that assertion.  The lower ideal may still be used by restricting to the full subcategory of objects which admit an arrow to $t_0$, but this is only cofinality on reachability preorders.  Comparison of total categories requires the actual comma categories.  The commentary formulates and proves the corresponding sufficient-and-necessary condition for base change to preserve the homotopy type of every Grothendieck fibration, using Quillen's Theorem~A \citep[\S1]{Quillen1973}.  The source's separate properness and smoothness criteria are also proved under its hypothesis on parallel pullbacks.  Contractibility of the strict fibres of $T\to T_{\mathrm o}$, together with local constancy of the coefficients there, allows homotopy Kan extension to descend the diagram to the preorder.  This argument explains why those sufficient criteria can hold without the universal comma condition \citep{Maltsiniotis2005Asphericite,Cisinski2009LocallyConstant}.

The sequential colimit itself is less problematic.  The nerve preserves filtered colimits, and for a sequence its homotopy colimit maps by a weak equivalence to its ordinary colimit.  We prove that two cofinal chains give isomorphic objects of $\Hot$ under the source's homotopy hypothesis: interleave cofinal subsequences, compare their bonding maps by homotopies, and compare the resulting telescopes.  This answers the existence part of the independence question.  The proof neither chooses a preferred isomorphism nor promotes it to a categorical homotopy equivalence in $\Cathot$.

Finite probes explain part of the remaining caution.  A finite category is finitely presented for ordinary categorical filtered colimits, which gives the asserted formula for $\Cathot$-morphisms into $C_\infty$.  Its nerve need not have the homotopy type of a finite CW complex.  The one-object category of the group of order two is a concrete counterexample to unrestricted compactness after passage to $\Hot$: its mod-two cohomology persists in every degree, whereas each finite-dimensional skeleton of its nerve has bounded cohomological dimension.  Replacing $\Cathot$ by $\Hot$ therefore needs a separate argument.  The mapping-set formulas at the end of the excerpt are audited with their parameter $Y$ and their ambient ind-category restored in the commentary.

The proposed obstruction from strict loop multiplication on manuscript page~113 is already withdrawn in footnote~24.  Kan's loop-group theorem supplies strict simplicial group models \citep[Theorem~7.1]{Kan1958Combinatorial}.  For unrestricted paths, the specific comparison is taken up positively later in this same chapter, \S13; it is not a current open problem.  The constant-path lemma left unwritten at manuscript page~170 is supplied here by successively retracting the end vertices of an oriented interval.  Modern hammock localization gives a separate, precise mapping-space theory.  Its requirement that vertical arrows also be weak equivalences must be checked when comparing it with Grothendieck's full path fibres \citep{DwyerKan1980Calculating,ArakawaCnossen2026}.
\nocite{*}
